

\section{Related Works}

\noindent\textbf{Regional dropout: }
Methods~\cite{devries2017cutout,zhong2017randomerase} removing random regions in images have been proposed to enhance the generalization performance of CNNs. 
Object localization methods~\cite{singh2017hide,choe2019attention} also utilize the regional dropout techniques for improving the localization ability of CNNs.
CutMix is similar to those methods, while the critical difference is that the removed regions are filled with patches from another training images.
DropBlock~\cite{ghiasi2018dropblock} has generalized the regional dropout to the feature space and have shown enhanced generalizability as well. 
CutMix can also be performed on the feature space, as we will see in the experiments.


\noindent\textbf{Synthesizing training data: }
Some works have explored synthesizing training data for further generalizability.
Generating new training samples by Stylizing ImageNet~\cite{ImageNet,geirhos2018imagenet_stylize} has guided the model to focus more on shape than texture, leading to better classification and object detection performances. 
CutMix also generates new samples by cutting and pasting patches within mini-batches, leading to performance boosts in many computer vision tasks; unlike stylization as in~\cite{geirhos2018imagenet_stylize}, CutMix incurs only negligible additional cost for training.
For object detection, object insertion methods~\cite{dwibedi2017cut_paste_learn, dvornik2018modeling} have been proposed as a way to synthesize objects in the background.
These methods aim to train a good represent of a single object samples, while CutMix generates combined samples which may contain multiple objects.

\noindent\textbf{Mixup: }
CutMix shares similarity with Mixup~\cite{zhang2017mixup, tokozume2018betweenclass} in that both combines two samples, where the ground truth label of the new sample is given by the linear interpolation of one-hot labels. 
As we will see in the experiments, Mixup samples suffer from the fact that they are locally ambiguous and {unnatural}, and therefore confuses the model, especially for localization.
Recently, Mixup variants~\cite{verma2018manifoldmixup, summers2019improvedm, guo2018adamix, takahashi2018ricap} have been proposed; they perform feature-level interpolation and other types of transformations.
Above works, however, generally lack a deep analysis in particular on the localization ability and transfer-learning performances.
We have verified the benefits of CutMix not only for an image classification task, but over a wide set of localization tasks and transfer learning experiments. 

\noindent\textbf{Tricks for training deep networks: }
Efficient training of deep networks is one of the most important problems in computer vision community, as they require great amount of compute and data.
Methods such as weight decay, dropout~\cite{Dropout}, and Batch Normalization~\cite{BN} are widely used to efficiently train deep networks. 
Recently, methods adding noises to the internal features of CNNs ~\cite{stochasticdepth,ghiasi2018dropblock,yamada2018shakedrop} or adding extra path to the architecture~\cite{SENet,GENet} have been proposed to enhance image classification performance. 
CutMix is complementary to the above methods because it operates on the data level, without changing internal representations or architecture.

